% !TeX root = ../教材题库_全部条目_转写工作稿.tex
% Source: XJ-B1-C05-RV-P207-Y1129
% 第5章 / 章末复习 复习题五 / 教材 p.207、208 / PDF 第214、215页
% OCR未校对：公式与文字均可能有误；未完成数学TeX转写；禁止直接发给学生
% 原题及相邻原解截图（核对用）：../题目原图/XJ-B1-C05-RV-P207-Y1129-01.webp, ../题目原图/XJ-B1-C05-RV-P207-Y1129-02.webp
% 本文件只在转写完成后作为题干使用；OCR 条目尚不是可靠题干。
\begin{bankitem}{复习题五 第15题}{第5章 / 章末复习 复习题五 / 教材 p.207、208 / PDF 第214、215页}{XJ-B1-C05-RV-P207-Y1129}{OCR未校对：公式与文字均可能有误；未完成数学TeX转写；禁止直接发给学生}
\begin{ocrdraft}
[教材p.207 / PDF第214页]
\par
15．如图都是函数3，一，Asin（-十@（，A>0，(D>()，》到<丌）在一个周期内的图
\par
象，试分别写出这两个函数的解析式．
\par
2
\par
2
\par
訌
\par
2是
\par
。4是
\par
是
\par
3
\par
3
\par
（2）
\par
（第巧题）
\par
[教材p.208 / PDF第215页]
\par
温故而知新
\par
1十sin
\end{ocrdraft}
\end{bankitem}
